\chapter{Nonlinear Models, Memory, and Stability Beyond the NEP}

The rational nonlinear eigenvalue problem developed in the previous part is
nonlinear in the spectral parameter \(s\). It is still a small-signal object:
it comes from a linearization in time and from exact or approximate
elimination of linearized states. This chapter separates that NEP theory from
the genuinely nonlinear power-system problem. The latter concerns equilibria,
local stability, bifurcations, regions of attraction, amplitude-dependent
damping, and memory generated by eliminating nonlinear controllers.

\thesisbox{thesisorange}{Separation of roles}{
The Schur NEP answers: what rational operator preserves the linearized
controller memory around an operating point? The nonlinear chapter answers:
which equilibria exist, how they lose stability, and whether a spectral margin
survives finite perturbations.}

\section{Nonlinear DAE Layer}

Let the full converter-rich model be written as a semi-explicit DAE,
\begin{align}
  \dot x &= f(x,y,p), \label{eq:nonlin_dae_f}\\
  0 &= g(x,y,p), \label{eq:nonlin_dae_g}
\end{align}
where \(x\) collects dynamic states, \(y\) collects algebraic network
variables, and \(p\) is an operating or planning parameter. A regular
equilibrium \((x^\star,y^\star)\) satisfies
\begin{equation}
  f(x^\star,y^\star,p)=0,\qquad
  g(x^\star,y^\star,p)=0,
  \qquad
  g_y(x^\star,y^\star,p)\ \hbox{nonsingular}.
  \label{eq:nonlin_equilibrium_regular}
\end{equation}
When the algebraic Jacobian is nonsingular, the implicit-function theorem gives
a local algebraic manifold \(y=h(x,p)\). The reduced ODE is
\begin{equation}
  \dot x=F(x,p)=f(x,h(x,p),p),
  \label{eq:nonlin_reduced_ode}
\end{equation}
and the small-signal state matrix is
\begin{equation}
  A_{\rm red}
  =
  f_x-f_y g_y^{-1}g_x
  \label{eq:nonlin_reduced_jacobian}
\end{equation}
at the equilibrium. The NEP chapters use this Jacobian, or a structured
second-order projection of it. The nonlinear theory must also study what
happens when the equilibrium moves, becomes singular, or ceases to attract
finite disturbances.

\section{Canonical Converter-Rich Family}

A useful hierarchical model for later chapters is
\begin{align}
  M\ddot\theta+D_0\dot\theta+P_e(\theta,V)+Gz &= p, \label{eq:nonlin_angle_layer}\\
  \dot z &= f_c(z,\theta,\dot\theta,V,p), \label{eq:nonlin_control_layer}\\
  0 &= g(\theta,V,z,p). \label{eq:nonlin_network_layer}
\end{align}
The state \(z\) may represent a synchronous-machine exciter or governor, a
grid-following PLL and current loop, a grid-forming droop or virtual
synchronous-machine controller, a voltage loop, a filter, or a plant-level
control state. Limiters and protections should be treated as smooth models
only inside a declared operating region; outside it they create hybrid
dynamics and require a different analysis.

Linearizing \eqref{eq:nonlin_angle_layer}--\eqref{eq:nonlin_network_layer}
recovers the QEP or Schur NEP used earlier. The nonlinear chapter keeps the
prelinearized equations visible so that a small-signal margin is not mistaken
for a large-signal certificate.

\section{Linear Controller Elimination as Memory}

For a linear controller block
\begin{equation}
  \dot z=A_cz+B_\theta\theta+B_\omega\dot\theta,
  \label{eq:nonlin_linear_controller}
\end{equation}
the exact variation-of-constants formula is
\begin{equation}
  z(t)
  =
  e^{A_ct}z(0)
  +
  \int_0^t
  e^{A_c(t-\tau)}
  \left[
  B_\theta\theta(\tau)+B_\omega\dot\theta(\tau)
  \right]d\tau .
  \label{eq:nonlin_linear_memory}
\end{equation}
Substitution into the retained equation produces an integro-differential
model. Its temporal kernel is
\begin{equation}
  K(t)=Ge^{A_ct}B,
  \label{eq:nonlin_memory_kernel}
\end{equation}
and its Laplace transform is the Schur-resolvent term
\begin{equation}
  G(sI-A_c)^{-1}B.
  \label{eq:nonlin_memory_transform}
\end{equation}
Thus the rational self-energy \(\Pi(s)\) is the spectral representation of
linear memory. This is the conceptual bridge between the NEP part and the
nonlinear-memory part of the dissertation.

\begin{theorem}[Self-energy as the tangent of memory]
\label{thm:nonlin_self_energy_tangent}
Consider a smooth nonlinear retained-hidden system near a regular equilibrium,
and suppose the hidden variables can be eliminated locally as a causal memory
functional acting on the retained trajectory. If this memory functional is
Fr\'echet differentiable at the equilibrium trajectory, then the Laplace
transform of its first variation is the Schur self-energy \(\Pi(s)\) of the
linearized hidden block.
\end{theorem}

\begin{proof}
Linearize the full nonlinear DAE or ODE around the equilibrium and partition
the linearization into retained and hidden coordinates. The hidden linear
subsystem has the variation-of-constants representation
\eqref{eq:nonlin_linear_memory}. Its contribution to the retained equation is
a convolution with kernel \(K(t)=Ge^{A_ct}B\). Taking the Laplace transform
gives \(G(sI-A_c)^{-1}B\), with the appropriate stiffness and velocity-channel
blocks collected as \(\Pi(s)\). This is exactly the first variation of the
nonlinear memory functional.
\end{proof}

\section{Nonlinear Elimination}

If the controller block is nonlinear,
\begin{equation}
  \dot z=f_c(z,\theta,\dot\theta,V,p),
  \label{eq:nonlin_nonlinear_controller}
\end{equation}
then eliminating \(z\) no longer gives a single matrix-valued rational function.
The reduced dynamics may contain history-dependent nonlinear functionals,
amplitude-dependent damping, and higher-order kernels. Locally, one can
organize those effects through:
\begin{itemize}
\item Volterra-series kernels when the nonlinear response is analytic and
small enough for a convergent expansion;
\item center-manifold or invariant-manifold reductions near critical
eigenvalues;
\item Mori-Zwanzig-type memory terms when unresolved variables are projected
out of the dynamics;
\item amplitude-dependent damping and stiffness coefficients when only the
dominant oscillatory envelope is retained.
\end{itemize}
The thesis will use this list as a research contract. A nonlinear memory claim
must state which reduction is being used and which finite-amplitude validation
gate supports it.

The same idea can be expressed as an amplitude-dependent self-energy when a
dominant oscillatory envelope is meaningful:
\begin{equation}
  \Pi(s;A)
  =
  \Pi_1(s)
  +
  A^2\Pi_3(s)
  +
  A^4\Pi_5(s)
  +\cdots .
  \label{eq:nonlin_amplitude_dependent_pi}
\end{equation}
Here \(\Pi_1(s)\) is the small-signal self-energy from Chapter~6, while higher
terms represent finite-amplitude stiffness and damping corrections. This
expansion is only legitimate under a declared Volterra, harmonic-balance, or
normal-form approximation. Its value is conceptual and testable: it asks
whether a small-signal damping reversal survives finite amplitude, disappears
through saturation, or becomes a Hopf mechanism. Recent non-Markovian closure
work reinforces the same principle that hidden variables generally reappear as
memory terms when only partial state information is retained
\cite{monsel2024neuraldde}.

\section{Bifurcation and Large-Signal Stability Contracts}

Small-signal stability asks whether the eigenvalues of the linearization have
negative real part. Nonlinear stability asks additional questions:
\begin{enumerate}
\item Does a regular equilibrium exist for the planned operating point?
\item Does the equilibrium lose stability through a Hopf, saddle-node, or
other bifurcation as \(p\) changes?
\item If a Hopf bifurcation occurs, what is the first Lyapunov coefficient and
does the emerging oscillation remain bounded?
\item What region of attraction is certified by an energy, Lyapunov,
contraction, or incremental-passivity argument?
\item Does a controller-induced damping reversal seen in the NEP persist under
finite perturbations, or is it only a local pole-motion effect?
\end{enumerate}

For a simple Hopf pair \(s(p)=\sigma(p)\pm j\omega(p)\), the spectral crossing
condition is
\begin{equation}
  \sigma(p_\star)=0,\qquad
  \omega(p_\star)\ne0,\qquad
  \frac{d\sigma}{dp}(p_\star)\ne0.
  \label{eq:nonlin_hopf_crossing}
\end{equation}
This condition locates a candidate bifurcation. It does not state the
amplitude, stability, or basin of the post-bifurcation motion. Those require
the nonlinear terms of \(F(x,p)\), not only the Schur NEP.

\section{Chapter Deliverables}

The nonlinear module is complete only when it contains the following
deliverables:
\begin{enumerate}
\item a regular DAE equilibrium theorem or a clearly stated regularity
assumption for each benchmark family;
\item a local-stability result tied to the QEP or Schur NEP;
\item at least one bifurcation calculation with finite-difference or complex
step verification of the derivative;
\item an energy, passivity, contraction, or Lyapunov-region argument for a
restricted model;
\item a finite-perturbation simulation showing whether the small-signal margin
predicts the observed nonlinear response;
\item a paper extraction section that states which nonlinear-memory result is
new, which result is known, and which result is only numerical evidence.
\end{enumerate}

This chapter is therefore not a replacement for the Schur NEP. It is the next
layer above it. The NEP preserves linearized memory; nonlinear analysis decides
which parts of that memory survive beyond infinitesimal perturbations.
